\documentclass[11pt]{amsart}
\usepackage{amsmath,amssymb,amsthm,booktabs}
\usepackage[margin=1.15in]{geometry}
\usepackage[hidelinks]{hyperref}
\newtheorem{theorem}{Theorem}[section]
\newtheorem{lemma}[theorem]{Lemma}
\newtheorem{proposition}[theorem]{Proposition}
\newtheorem{corollary}[theorem]{Corollary}
\newtheorem{conjecture}{Conjecture}
\theoremstyle{remark}
\newtheorem{remark}[theorem]{Remark}
\newtheorem{example}[theorem]{Example}
\newcommand{\PP}{\mathbb{P}}
\newcommand{\E}{\mathbb{E}}
\newcommand{\Z}{\mathbb{Z}}
\newcommand{\doi}[1]{\href{https://doi.org/#1}{doi:\nolinkurl{#1}}}
\newcommand{\arxiv}[1]{\href{https://arxiv.org/abs/#1}{arXiv:\nolinkurl{#1}}}
\setlength{\emergencystretch}{2.5em}

\title[FA-1f out of equilibrium at high density]{Relaxation of the one-dimensional Fredrickson--Andersen model from
a single infection at high density}
\author{Yitzchak Shmalo}
\address{Einstein Institute of Mathematics, The Hebrew University of Jerusalem, Jerusalem, Israel}
\email{yspennstate@gmail.com}
\date{November 2026}
\subjclass[2020]{60K35, 82C22}
\keywords{kinetically constrained models, Fredrickson--Andersen model, relaxation to equilibrium, out of equilibrium}

\raggedbottom
\begin{document}
\begin{abstract}
Consider the one-spin facilitated Fredrickson--Andersen model on $\Z$ with equilibrium density $q$ of infected
sites. For $q>1/2$ we prove that, for every initial law and every local function, the difference between the expectation
at time $t$ and the equilibrium mean is bounded by an exponentially small term plus the probability that the
initial distance from the origin to the nearest infection exceeds a fixed multiple of $t$; the constants depend
only on $q$ and the support of the function. No independence, stationarity
or uniform control of the gaps between infections is assumed. Hence the two conjectures of Martinelli, Shapira and
Toninelli on convergence from any law with an infection, and on exponential relaxation when the distance from the
origin to the nearest infection has an exponential tail, hold for every $q>1/2$; this includes a single initial
infection. The proof transfers the relaxation theorem of Blondel, Cancrini, Martinelli, Roberto and Toninelli from
uniform to origin-centred hypotheses.
\end{abstract}
\maketitle

\section{Introduction}\label{sec:intro}

Let $\Omega=\{0,1\}^{\Z}$. A site $x$ with $\eta_x=0$ is \emph{infected} and a site with $\eta_x=1$ is healthy. Fix
$q\in(0,1)$, put $p=1-q$, and let $\pi$ be the product measure on $\Omega$ with $\pi(\eta_x=0)=q$ for every $x$. The
one-spin facilitated Fredrickson--Andersen model (FA-1f) \cite{FA84} is the Markov process on $\Omega$ with generator
\begin{equation}\label{eq:generator}
Lf(\eta)=\sum_{x\in\Z}c_x(\eta)\bigl(\pi_x(f)-f\bigr)(\eta),\qquad c_x(\eta)=1-\eta_{x-1}\eta_{x+1},
\end{equation}
where $\pi_x(f)$ averages $f$ over the spin at $x$ with respect to $\pi$. Thus every site carries a rate-one clock,
and at a ring the spin at $x$ is replaced by a fresh spin, infected with probability $q$, provided that at least one
neighbour of $x$ is infected. The process is reversible with respect to $\pi$. We write $P_t$ for its semigroup and
$\nu P_t f=\int P_tf\,d\nu$ for a probability measure $\nu$ on $\Omega$, called the \emph{initial law}. A function
$f$ on $\Omega$ is local if it depends on finitely many spins; its support is the smallest such set of sites. For
$x\in\Z$ let
\begin{equation}\label{eq:dist}
D_x(\eta)=\min\{|x-y|:\eta_y=0\},\qquad R(\eta)=D_0(\eta),
\end{equation}
with $\min\emptyset=\infty$. The following two conjectures are stated, in equivalent form, in the survey of Martinelli, Shapira and
Toninelli \cite[Section~1.1, Conjectures~1 and~2]{MST25}; the first is also Conjecture~2.11 in \cite{Mar26}.

\begin{conjecture}\label{conj:convergence}
For every $q\in(0,1)$, every initial law $\nu$ with $\nu(R<\infty)=1$ and every local function $f$,
$\lim_{t\to\infty}\nu P_tf=\pi(f)$.
\end{conjecture}

\begin{conjecture}\label{conj:rate}
Let $q\in(0,1)$ and let $\nu$ satisfy $\nu(R>\ell)\le Ae^{-\kappa\ell}$ for all $\ell\ge0$, with $A<\infty$ and
$\kappa>0$. Then for every local function $f$ there are $C,c>0$ with $|\nu P_tf-\pi(f)|\le Ce^{-ct}$ for all
$t\ge0$.
\end{conjecture}

The condition $\nu(R<\infty)=1$ says that the initial configuration has at least one infection almost surely; it
cannot be dropped, because the configuration without infections does not move. The difficulty is that FA-1f is not
attractive and its constraint degenerates, so that neither monotone coupling nor the positive spectral gap of the
equilibrium dynamics \cite{CMRT08} applies to an initial law that is singular with respect to $\pi$.

Several cases are known. Blondel, Cancrini, Martinelli, Roberto and Toninelli \cite{BCMRT13} proved exponential
relaxation for $q>1/2$ under the uniform hypothesis $\sup_x\nu(\theta^{D_x})<\infty$ for some $\theta>1$, which holds
for every nontrivial product law; it fails for a single infection, where $D_x(\eta)=|x|$. Mountford and Valle
\cite{MV19} proved exponential relaxation from a single infection, and from origin-centred initial laws as in
Conjecture~\ref{conj:rate}, for $q$ sufficiently close to $1$. Cancrini, Martinelli, Schonmann and Toninelli
\cite{CMST10} proved exponential relaxation for every $q$ when the finite-window densities of $\nu$ with respect to
$\pi$ grow at most like a small exponential of the window size. The fronts of FA-1f were studied by Blondel,
Deshayes and Toninelli \cite{BDT19}, and \cite[Theorem~6.2]{MST25} shows that a finite set of infections spreads
linearly for every $q$.

The main result removes the uniformity from the hypothesis of \cite{BCMRT13}.

\begin{theorem}\label{thm:main}
Let $q\in(1/2,1)$ and let $S\subset\Z$ be finite. There are $\gamma,c,C>0$ depending only on $q$ and $S$ such that
for every initial law $\nu$, every function $f$ with support in $S$ and every $t\ge0$,
\begin{equation}\label{eq:main}
|\nu P_tf-\pi(f)|\le\|f\|_\infty\bigl(Ce^{-ct}+2\,\nu(R>\gamma t)\bigr).
\end{equation}
\end{theorem}

\begin{corollary}\label{cor:conjectures}
Let $q\in(1/2,1)$. Conjectures~\ref{conj:convergence} and~\ref{conj:rate} hold. More precisely, under the
hypothesis of Conjecture~\ref{conj:rate} the rate can be taken as $\min\{c,\kappa\gamma\}$ with $c,\gamma$ from
Theorem~\ref{thm:main}, and if $\nu(R>\ell)\le Ae^{-\kappa\ell^\beta}$ with $0<\beta<1$, then
$|\nu P_tf-\pi(f)|\le C'e^{-c't^\beta}$. In particular, from a single infection at any site, every local function
converges to its equilibrium mean exponentially fast.
\end{corollary}

Theorem~\ref{thm:main} has no uniform rate over all initial laws, and none is possible: if the only infection sits
at distance $\ell$ from the support of $f$, the spins on the support cannot change before a chain of clock rings has
crossed the distance $\ell$, which takes time of order $\ell$ (Lemma~\ref{lem:fp}). The term $\nu(R>\gamma t)$
accounts for this.

The proof has two parts. A Lyapunov estimate (Lemma~\ref{lem:drift}) shows that when $q>1/2$ the exponential moment
$\E\,\theta^{D_x}$ at a fixed site forgets its initial value at an exponential rate. Running the process for most of
the time interval, it follows that at sites within distance of order $t$ from the origin these moments are bounded
uniformly. We then replace the spins far from the origin by infections, which makes the moment bound uniform over
all sites, apply the theorem of \cite{BCMRT13} on the remaining time interval, and remove the replacement by finite
speed of propagation. The restriction $q>1/2$ enters only through Lemma~\ref{lem:drift}; see
Remark~\ref{rem:half}.

Section~\ref{sec:prelim} collects the graphical construction, finite speed of propagation and the input from
\cite{BCMRT13}. Lemma~\ref{lem:drift} is proved in Section~\ref{sec:drift}, and Theorem~\ref{thm:main} and
Corollary~\ref{cor:conjectures} in Section~\ref{sec:proof}. Section~\ref{sec:further} contains two complements: the
two extreme infections of a finite infected set move away from each other at linear speed when $q>1/2$
(Proposition~\ref{prop:fronts}), and, for every $q$, the argument of \cite{CMST10} with explicit constants, applied
to a class of correlated initial laws (Proposition~\ref{prop:allq}). A table of notation is in
Appendix~\ref{app:notation}.

\section{Preliminaries}\label{sec:prelim}

This section fixes the construction of the process, proves finite speed of propagation in the form needed below,
and states the relaxation theorem of \cite{BCMRT13}. Throughout, $C,C',C_0$ and $c,c',c_0$ denote positive
constants whose dependence is stated where they are introduced; they do not change from line to line.

\subsection{Graphical construction}
Attach to each site $x\in\Z$ a Poisson process of rate one, the clock of $x$, and to its $k$-th ring an independent
spin $s_{x,k}\in\{0,1\}$ with $\PP(s_{x,k}=0)=q$, all clocks and spins independent. Given $\eta\in\Omega$, the
process $(\eta(t))_{t\ge0}$ with $\eta(0)=\eta$ is defined by setting $\eta_x(t)=s_{x,k}$ at the $k$-th ring of the
clock of $x$ if $c_x(\eta(t-))=1$, and leaving the configuration unchanged otherwise. Since the constraint has range
one, this defines the process on $\Z$ almost surely for every initial configuration, and its generator is
\eqref{eq:generator}; see \cite{CMRT08}. We write $\PP_\eta$ and $\E_\eta$ for the law and expectation of this
process, so that $P_tf(\eta)=\E_\eta f(\eta(t))$.

\subsection{Finite speed of propagation}
The following lemma is standard; we include the proof because it is used for processes whose dynamics outside a
finite interval differ from FA-1f.

\begin{lemma}\label{lem:fp}
Let $I\subset\Z$ be an interval and $S\subset I$ finite, and let $d\ge1$ be the distance from $S$ to $\Z\setminus I$.
Let $\eta(\cdot)$ and $\zeta(\cdot)$ be two processes built from the same clocks and spins, which at every site of
$I$ follow the update rule of FA-1f, with arbitrary (possibly different) evolution outside $I$, adapted to the
clocks and spins, and with $\eta_x(0)=\zeta_x(0)$ for $x\in I$. Then for $u\ge0$,
\begin{equation}\label{eq:fp}
\PP\bigl(\eta_y(\tau)\ne\zeta_y(\tau)\ \text{for some }y\in S\text{ and some time }\tau\le u\bigr)\le|S|\,\frac{(2u)^d}{d!}.
\end{equation}
If $d\ge8eu$, the right-hand side is at most $|S|\,4^{-d}$.
\end{lemma}

\begin{proof}
Suppose $\eta_y(\tau)\ne\zeta_y(\tau)$ with $y\in I$ and $\tau\le u$. The spins at $y$ agreed at time $0$ and change only at
rings of the clock of $y$, so there is a last ring of that clock at a time $\tau_1\le\tau$ just before which the spins at
$y$ agreed and after which they disagree. Both processes use the spin $s_{y,k}$ of that ring, so the constraints
$c_y$ of the two configurations differed just before $\tau_1$; hence a neighbour $y_1$ of $y$ had disagreeing spins
at time $\tau_1-$. If $y_1\in I$, repeat the argument from $(y_1,\tau_1-)$. This produces sites $y_0=y,y_1,y_2,\dots$, each a
neighbour of the previous one, and times $\tau\ge\tau_1>\tau_2>\dots$ such that the clock of $y_{j-1}$ rings at
$\tau_j$, and
the construction stops at the first $y_j\notin I$. If $y\in S$, at least $d$ steps are needed to leave $I$, so the
clocks of $y_0,\dots,y_{d-1}$ ring at times $u\ge\tau_1>\dots>\tau_d\ge0$. For a fixed sequence $y_0,\dots,y_{d-1}$ of
sites, the expected number of such decreasing $d$-tuples of rings is $u^d/d!$, by the factorial moment formula of a
Poisson process (also when sites repeat), and there are at most $|S|\,2^{d-1}$ nearest-neighbour sequences starting
in $S$. Markov's inequality gives \eqref{eq:fp}. If $d\ge8eu$, then $(2u)^d/d!\le(2eu/d)^d\le4^{-d}$, using
$d!\ge(d/e)^d$.
\end{proof}

\subsection{Relaxation under a uniform moment bound}
The input from \cite{BCMRT13} is the following.

\begin{theorem}[{\cite[Theorem~2.1]{BCMRT13}}]\label{thm:bcmrt}
Let $q\in(1/2,1)$, $\theta>1$, $K<\infty$ and $S\subset\Z$ finite. There are $C_0,c_0>0$ depending only on
$q,\theta,K$ and $S$ such that for every initial law $\mu$ with $\sup_{x\in\Z}\mu(\theta^{D_x})\le K$, every $f$
with support in $S$ and every $h\ge0$,
\begin{equation}\label{eq:bcmrt}
|\mu P_hf-\pi(f)|\le C_0\|f\|_\infty e^{-c_0h}.
\end{equation}
\end{theorem}

\begin{remark}\label{rem:bcmrt}
The initial law in \cite{BCMRT13} is arbitrary; independence is not assumed. The statement there lists the
dependence of the constants on $q$, the support and the moment bound, and the proof shows that $\mu$ enters only
through $\theta$ and $K$: after a reduction to an interval of length of order $h$, divided into blocks of length
proportional to $h$, the proof compares the process with the process restricted to configurations with at least
two infections in every block, and $\mu$ is used only to bound $\E_\mu\theta^{D_x(\eta(\tau))}$ at the block centres for
$\tau\le h$, through the drift estimate of \cite[Proposition~4.1]{BCMRT13}, followed by a union bound over the blocks. In
\cite[Section~3.2]{BCMRT13} the probability that the process leaves the set of good configurations before time $h$
should be bounded by $(1+|\Lambda|h)\sup_{\tau\le h}\PP(\eta(\tau)\ \text{is not good})$, where $|\Lambda|$ is the size of
the interval, as in \cite[Section~3]{BDT19}. Since $|\Lambda|$ grows linearly in $h$, this factor changes only the
constants. In the proof below $\theta$ and $K$ are fixed in terms of $q$ alone.
\end{remark}

\section{The distance to the nearest infection}\label{sec:drift}

For $q\in(1/2,1)$ we fix throughout
\begin{equation}\label{eq:constants}
\theta=\sqrt{q/p},\qquad a=\log\theta,\qquad \lambda=(\theta-1)\Bigl(\frac q\theta-p\Bigr),\qquad
B=\frac{q}{q-p\theta}.
\end{equation}
Since $1<\theta<q/p$, both $\lambda$ and $B$ are positive.

\begin{lemma}\label{lem:drift}
Let $q\in(1/2,1)$. For every $\eta\in\Omega$, every $x\in\Z$ with $D_x(\eta)<\infty$ and every $s\ge0$,
\begin{equation}\label{eq:drift}
\E_\eta\bigl[\theta^{D_x(\eta(s))}\bigr]\le e^{-\lambda s}\,\theta^{D_x(\eta)}+B.
\end{equation}
\end{lemma}

\begin{proof}
We first bound the generator applied to $g=\theta^{D_x}$. Let $D_x(\eta)=d\ge1$ and let $y$ be an infected site with
$|x-y|=d$. The neighbour $y'$ of $y$ between $x$ and $y$ is healthy, is at distance $d-1$ from $x$, and satisfies
$c_{y'}(\eta)=1$; it becomes infected at rate $q$, and then $D_x$ drops by one. A transition increases $D_x$ only if
an infected site at distance $d$ from $x$ heals while it is the only infected site at that distance; this happens
at rate at most $p$, and since a site can heal only next to an infected neighbour, which here is at distance $d+1$
from $x$, the distance becomes $d+1$. All other transitions do not increase $D_x$. Hence
\[
Lg(\eta)\le q\,(\theta^{d-1}-\theta^d)+p\,(\theta^{d+1}-\theta^d)=-\lambda\,\theta^d.
\]
If $d=0$, the site $x$ heals at rate at most $p$, and then $D_x=1$, so $Lg(\eta)\le p(\theta-1)$. Since
$-\lambda+q(\theta-1)/\theta=p(\theta-1)$, in both cases
\begin{equation}\label{eq:lyapunov}
Lg\le-\lambda g+b,\qquad b=\frac{q(\theta-1)}{\theta}.
\end{equation}
The same bound, with the same proof, holds for the process on an interval $\Lambda_N=[-N,N]$ in which the sites
outside $\Lambda_N$ are infected and never updated, and in which $D_x$ counts these sites as infected: the frozen
sites cannot heal, and a healthy site next to an infected one, frozen or not, still becomes infected at rate $q$.
This process has finitely many states, so Dynkin's formula applies to it, and \eqref{eq:lyapunov} gives
$\frac{d}{ds}\E g\le-\lambda\E g+b$, whence $\E g(s)\le e^{-\lambda s}g(0)+b/\lambda$. Note that $b/\lambda=B$.
Start this process from the configuration $\eta^N$ equal to $\eta$ on $\Lambda_N$ and infected outside; then
$D_x(\eta^N)\le D_x(\eta)$.

Fix $M\ge1$ and an integer $k$ with $\theta^k\ge M$. The function $\min\{\theta^{D_x},M\}$ depends only on the spins
in $[x-k,x+k]$. Build the process on $\Lambda_N$ and the process on $\Z$ from the same clocks and spins. Lemma~\ref{lem:fp}
with $I=\Lambda_N$ and $S=[x-k,x+k]$ shows that the two agree on $S$ at time $s$ with probability tending to one as
$N\to\infty$. Therefore
\[
\E_\eta\min\{\theta^{D_x(\eta(s))},M\}=\lim_{N\to\infty}\E\min\{\theta^{D_x(\eta^N(s))},M\}
\le e^{-\lambda s}\theta^{D_x(\eta)}+B,
\]
and letting $M\to\infty$ gives \eqref{eq:drift} by monotone convergence.
\end{proof}

The bound \eqref{eq:lyapunov} is the infinite-volume form of \cite[Proposition~4.1]{BCMRT13}. For $D_x\ge1$ the estimate
$Lg\le-\lambda g$ is sharp: equality holds when $x$ has a single nearest infection $y$, the site $y$ has an infected neighbour on the side away
from $x$, and no other transition changes $D_x$.

\section{Proof of Theorem~\ref{thm:main} and Corollary~\ref{cor:conjectures}}\label{sec:proof}

Fix $q\in(1/2,1)$, a finite $S\subset\Z$ and an integer $m\ge0$ with $S\subset[-m,m]$, and take
$\theta,a,\lambda,B$ from \eqref{eq:constants}. Set
\begin{equation}\label{eq:parameters}
\gamma=\frac{\lambda}{8a},\qquad \varepsilon=\min\Bigl\{\frac14,\frac{\gamma}{32e}\Bigr\},\qquad
h=\varepsilon t,\qquad s=(1-\varepsilon)t,\qquad r=\lfloor\gamma t\rfloor,
\end{equation}
and $t_0=2(m+1)/\gamma$. Let $C_0,c_0$ be the constants of Theorem~\ref{thm:bcmrt} for
$q$, $\theta$, $K=1+B$ and $S$; all of these depend only on $q$ and $S$.

\begin{proof}[Proof of Theorem~\ref{thm:main}]
\emph{Step 1: a moment bound near the origin.} Let $t\ge t_0$ and let $\eta$ be a configuration with
$R(\eta)\le\gamma t$. For $|x|\le r$ the triangle inequality gives $D_x(\eta)\le|x|+R(\eta)\le2\gamma t$, so by
Lemma~\ref{lem:drift} and $2a\gamma=\lambda/4$, $\varepsilon\le1/4$,
\begin{equation}\label{eq:step1}
\E_\eta\bigl[\theta^{D_x(\eta(s))}\bigr]\le\exp\Bigl(-\lambda(1-\varepsilon)t+\frac{\lambda t}{4}\Bigr)+B\le1+B
\qquad(|x|\le r).
\end{equation}

\emph{Step 2: replacement.} Define the random configuration $\zeta$ by $\zeta_x=\eta_x(s)$ for $|x|\le r$ and
$\zeta_x=0$ for $|x|>r$, and let $\mu$ be its law. For $|x|\le r$ we have $D_x(\zeta)\le D_x(\eta(s))$, and for
$|x|>r$ we have $D_x(\zeta)=0$. By \eqref{eq:step1},
\begin{equation}\label{eq:step2}
\sup_{x\in\Z}\mu\bigl(\theta^{D_x}\bigr)\le1+B.
\end{equation}
The law $\mu$ depends on $\eta$ and $t$, is correlated and is not stationary; Theorem~\ref{thm:bcmrt} allows this, and
its constants are the same for all these laws. Theorem~\ref{thm:bcmrt} gives
\begin{equation}\label{eq:step3}
|\mu P_hf-\pi(f)|\le C_0\|f\|_\infty e^{-c_0h}.
\end{equation}

\emph{Step 3: removing the replacement.} By the Markov property at time $s$, $P_tf(\eta)=\E_\eta[P_hf(\eta(s))]$,
while $\mu P_hf=\E_\eta[P_hf(\zeta)]$. Given $\eta(s)$, run two copies of FA-1f for time $h$ from $\eta(s)$ and from
$\zeta$ with the same fresh clocks and spins; their initial configurations agree on $[-r,r]$. Since $t\ge t_0$,
\[
r-m\ge\gamma t-1-m\ge\frac{\gamma t}{2}\ge8e\varepsilon t=8eh,
\]
so Lemma~\ref{lem:fp}, with $I=[-r,r]$ and $d\ge r-m$, gives
\begin{equation}\label{eq:step4}
\bigl|P_tf(\eta)-\mu P_hf\bigr|\le2\|f\|_\infty|S|\,4^{-\gamma t/2}.
\end{equation}

\emph{Step 4: conclusion.} By \eqref{eq:step3} and \eqref{eq:step4}, for $t\ge t_0$ and every $\eta$ with
$R(\eta)\le\gamma t$,
\[
|P_tf(\eta)-\pi(f)|\le\|f\|_\infty\bigl(C_0e^{-c_0\varepsilon t}+2|S|e^{-(\gamma\log2)t}\bigr).
\]
Put $c=\min\{c_0\varepsilon,\gamma\log2\}$ and $C=\max\{C_0+2|S|,\,2e^{ct_0}\}$. Since $|P_tf-\pi(f)|\le2\|f\|_\infty$
always, we obtain for all $t\ge0$
\begin{equation}\label{eq:pointwise}
\sup_{\eta:\,R(\eta)\le\gamma t}|P_tf(\eta)-\pi(f)|\le C\|f\|_\infty e^{-ct}.
\end{equation}
Integrating \eqref{eq:pointwise} against $\nu$ on $\{R\le\gamma t\}$ and using $|P_tf-\pi(f)|\le2\|f\|_\infty$ on
$\{R>\gamma t\}$ gives \eqref{eq:main}.
\end{proof}

\begin{proof}[Proof of Corollary~\ref{cor:conjectures}]
If $\nu(R<\infty)=1$, then $\nu(R>\gamma t)\to0$ as $t\to\infty$, and \eqref{eq:main} gives
Conjecture~\ref{conj:convergence}. Under the hypothesis of Conjecture~\ref{conj:rate},
$\nu(R>\gamma t)\le Ae^{-\kappa\gamma t}$, and \eqref{eq:main} gives the rate $\min\{c,\kappa\gamma\}$. The
stretched-exponential case is the same. For a single infection at a site $z$, $R=|z|$ is deterministic and the
second term of \eqref{eq:main} vanishes for $t\ge|z|/\gamma$.
\end{proof}

\section{Two complements}\label{sec:further}

The first result concerns the extreme infections of a finite infected set when $q>1/2$; the second gives exponential
relaxation for every $q$ from initial laws whose local densities with respect to $\pi$ are not too large.

\subsection{The extreme infections}
For a configuration with finitely many infections let $X_+(\eta)$ and $X_-(\eta)$ be the largest and the smallest
infected site. An infected site can heal only if it has an infected neighbour, so the set of infections never becomes
empty; by Lemma~\ref{lem:fp} it stays finite at all times.

\begin{proposition}\label{prop:fronts}
Let $q\in(1/2,1)$ and let $\eta$ have finitely many infections, at least one. Then $\PP_\eta$-almost surely
\[
2q-1\le\liminf_{t\to\infty}\frac{X_+(\eta(t))}{t}\le\limsup_{t\to\infty}\frac{X_+(\eta(t))}{t}\le q,
\]
and symmetrically $-q\le\liminf_t X_-(\eta(t))/t\le\limsup_t X_-(\eta(t))/t\le-(2q-1)$.
\end{proposition}

\begin{proof}
Write $X(t)=X_+(\eta(t))$. The site $X+1$ is healthy with an infected neighbour, so $X$ increases by one at rate $q$.
The site $X$ can heal only if $X-1$ is infected, and then $X$ decreases by one; this happens at rate $p$ if
$\eta_{X-1}=0$ and at rate $0$ otherwise. No other transition moves $X$. Hence
\[
X(t)-X(0)=M(t)+\int_0^t\bigl(q-p\,\mathbf 1\{\eta_{X(\tau)-1}(\tau)=0\}\bigr)\,d\tau,
\]
where $\mathbf 1$ denotes an indicator and $M$ is a martingale with jumps $\pm1$ whose predictable quadratic variation is at most $t$. By Doob's
inequality, $\PP(\sup_{\tau\le2^{n+1}}|M(\tau)|>\delta2^n)\le4\cdot2^{n+1}/(\delta^24^n)$, which is summable in $n$, so
$M(t)/t\to0$ almost surely by the Borel--Cantelli lemma. The integrand lies in $[q-p,q]=[2q-1,q]$. The statement for
$X_-$ follows by reflection.
\end{proof}

The diameter $X_+-X_-$ grows linearly for every $q\in(0,1)$ by \cite[Theorem~6.2]{MST25}. For $q>1/2$,
Proposition~\ref{prop:fronts} shows in addition that each extreme infection moves away from the origin at a linear
speed, and Corollary~\ref{cor:conjectures} shows that in every fixed finite window the law of the process converges to
equilibrium.

\subsection{Initial laws close to equilibrium, every q}
Let $g(q)>0$ be the spectral gap of FA-1f on $\Z$ in $L^2(\pi)$, which is positive for every $q\in(0,1)$
\cite{CMRT08}. The following is \cite[Theorem~2.1]{CMST10} for FA-1f, with explicit constants.

\begin{proposition}\label{prop:allq}
Let $q\in(0,1)$, $v=16e$, and let $\nu$ satisfy, for some $K<\infty$ and $0\le b<2g(q)/v$,
\begin{equation}\label{eq:density}
\max_{\omega\in\{0,1\}^{[-r,r]}}\frac{\nu(\eta_{[-r,r]}=\omega)}{\pi(\eta_{[-r,r]}=\omega)}\le Ke^{br}
\qquad(r=0,1,2,\dots).
\end{equation}
Then for every local $f$ there is $C<\infty$ with $|\nu P_tf-\pi(f)|\le Ce^{-c't}$ for all $t\ge0$, where
$c'=\min\{g(q)-bv/2,\ 16e\log4\}>0$.
\end{proposition}

\begin{proof}
Let $S\subset[-m,m]$ be the support of $f$ and $r=\lceil vt\rceil+m+1$. Let $\mu_r$ be the law that agrees with $\nu$
on $[-r,r]$ and is $\pi$ outside, independently. Lemma~\ref{lem:fp} with $d\ge vt\ge8et$ gives
$|\nu P_tf-\mu_rP_tf|\le2\|f\|_\infty|S|4^{-vt}$. The density $\phi_r=d\mu_r/d\pi$ depends only on the spins in
$[-r,r]$, and by \eqref{eq:density}, $\|\phi_r-1\|_{L^2(\pi)}^2=\pi(\phi_r^2)-1\le Ke^{br}$. Since $\pi$ is invariant and
$P_t$ contracts by $e^{-g(q)t}$ on functions of mean zero in $L^2(\pi)$,
\[
|\mu_rP_tf-\pi(f)|=\bigl|\pi\bigl((\phi_r-1)\,P_t(f-\pi(f))\bigr)\bigr|\le\sqrt K\,e^{br/2}e^{-g(q)t}\,2\|f\|_\infty,
\]
and $e^{br/2}\le e^{b(m+2)/2}e^{bvt/2}$.
\end{proof}

\begin{example}\label{ex:markov}
Let $\nu$ be the law of a stationary Markov chain $(\eta_x)_{x\in\Z}$ on $\{0,1\}$ with transition matrix $T$ having
positive entries, stationary vector $\alpha$, and $T_{ij}\le e^{\delta}\pi(\eta_0=j)$ for all $i,j$, where
$0\le\delta<g(q)/v$. The probability of a word of length $2r+1$ divided by its $\pi$-probability is at most
$\max_i\alpha_i/\pi(\eta_0=i)\cdot e^{2\delta r}$, so \eqref{eq:density} holds with $b=2\delta$, and
Proposition~\ref{prop:allq} gives exponential relaxation for every $q$. If the two rows of $T$ differ, $\nu$ is not
a product measure. Such laws satisfy the hypothesis of Conjecture~\ref{conj:rate}, since
$\nu(R>r)=\alpha_1T_{11}^{2r}$, but the condition on $\delta$ is much stronger than that hypothesis when $g(q)$ is
small.
\end{example}

\begin{remark}\label{rem:half}
Lemma~\ref{lem:drift} needs $\theta<q/p$, which is possible only for $q>1/2$. For $q\le1/2$, $\theta>1$ and $d\ge1$, take the
configuration whose infections are exactly the sites $d$ and $d+1$; for $g=\theta^{D_0}$ it has
$Lg=\theta^d(\theta-1)(p-q/\theta)\ge0$, so no bound of the form \eqref{eq:lyapunov} with $\lambda>0$ holds for
all large $d$. This is a statement about one
Lyapunov function, not about the process. For $q\le1/2$ both conjectures remain open, including the case of a single
infection. What is missing is a lower bound on the density of infections in a region of size of order $t$ around the
initial infection, strong enough to run the block argument of \cite{BCMRT13}.
\end{remark}

\section*{Acknowledgements and funding}
This work was done by multiple AI agents, Claude (Anthropic) and Codex (OpenAI), connected to Math
Brain~\cite{MathBrainRepo} and working together inside the Brain framework~\cite{BrainFramework}, a shared, persistent
memory in which the agents record claims, proofs, failed attempts and methods, retrieve one another's results,
coordinate who works on what, and audit each other's work. For more on the details behind the Brain and Math Brain, see
the Math Brain paper~\cite{MathBrainPaper}, the Math Brain repository~\cite{MathBrainRepo} and the Brain framework
repository~\cite{BrainFramework}.

The research presented in this paper was supported by the European Research Council (ERC) under the European Union's
Horizon Europe research and innovation programme (grant agreement No.~101041711), by the Simons Foundation as part of
the Collaboration on the Mathematical and Scientific Foundations of Deep Learning, by Heights Labs, by Convex Nexus
Capital, by the Israel Science Foundation (grant number 2258/19), and by the Israel Science Foundation (ISF Grant
4101/25). Computational resources were provided in part by Google Cloud through its education credits programme (credit
530550103).

\section*{Declaration of competing interest}
The author holds equity in Heights Labs and in Convex Nexus Capital and has a management role at Convex Nexus Capital.
These entities are included in the funding declaration above.

\section*{Declaration of generative AI and AI-assisted technologies}
Codex agents (OpenAI) proved Theorem~\ref{thm:main} and Propositions~\ref{prop:fronts} and~\ref{prop:allq}; two further
Codex agents checked the proofs as referees; Claude agents (Anthropic) checked the steps and prepared the manuscript.
The author directed the work and is responsible for the content of the paper and for decisions concerning its
dissemination.

\section*{Code availability}
A script that checks the generator inequality \eqref{eq:lyapunov}, in floating-point arithmetic, on every configuration of
windows of up to twelve sites with frozen infected exterior, for several values of $q>1/2$, and that simulates the process from a single
infection, is in the repository
\begin{center}\url{https://github.com/yspennstate/fa1f-relaxation-high-density/tree/v1.1}.\end{center}
In the simulations ($q=0.6$ and $q=0.75$, $300$ paths each) the probabilities that the origin is infected and that
the sites $-1,0,1$ are all healthy are within two standard errors of $q$ and $p^3$ at $30$ of the $32$ recorded points; the two
exceptions, the origin at $q=0.6$, $t=45$ and at $q=0.75$, $t=70$, are $2.3$ standard errors away.

\appendix
\section{Notation}\label{app:notation}

\begin{center}
\begin{tabular}{p{3.5cm}p{7.3cm}p{3.3cm}}
\toprule
Symbol & Meaning & First appears\\
\midrule
$\Omega$, $\eta$, $\eta_x$ & configurations, a configuration, the spin at $x$ ($0$ infected, $1$ healthy) & Section~\ref{sec:intro}\\
$q$, $p$ & equilibrium probability of infection, $p=1-q$ & Section~\ref{sec:intro}\\
$\pi$, $\pi_x$ & equilibrium product measure, average over the spin at $x$ & \eqref{eq:generator}\\
$c_x$, $L$, $P_t$ & constraint, generator, semigroup & \eqref{eq:generator}\\
$\nu$, $\mu$ & initial laws & Section~\ref{sec:intro}, Theorem~\ref{thm:bcmrt}\\
$f$, $S$ & local function, a finite set containing its support & Section~\ref{sec:intro}\\
$D_x$, $R$ & distance from $x$, and from $0$, to the nearest infection & \eqref{eq:dist}\\
$\PP_\eta$, $\E_\eta$ & law and expectation of the process started at $\eta$ & Section~\ref{sec:prelim}\\
$I$, $d$ & interval of agreement and distance from $S$ to its complement & Lemma~\ref{lem:fp}\\
$\theta$, $a$, $\lambda$, $B$ & constants of the drift estimate & \eqref{eq:constants}\\
$b$ & additive constant in the drift estimate & \eqref{eq:lyapunov}\\
$\Lambda_N$, $\eta^N$ & interval with frozen infected exterior, restricted configuration & Lemma~\ref{lem:drift}\\
$\gamma$, $\varepsilon$, $h$, $s$, $r$, $t_0$, $m$ & parameters of the proof of Theorem~\ref{thm:main} & \eqref{eq:parameters}\\
$\zeta$ & configuration with exterior replaced by infections & Section~\ref{sec:proof}\\
$C_0$, $c_0$ & constants of Theorem~\ref{thm:bcmrt} & \eqref{eq:bcmrt}\\
$X_\pm$ & largest and smallest infected site & Proposition~\ref{prop:fronts}\\
$g(q)$, $v$, $\phi_r$ & spectral gap, speed constant $16e$, density of $\mu_r$ & Proposition~\ref{prop:allq}\\
\bottomrule
\end{tabular}
\end{center}

\begin{thebibliography}{99}
\bibitem{BCMRT13} O.~Blondel, N.~Cancrini, F.~Martinelli, C.~Roberto and C.~Toninelli, Fredrickson--Andersen one spin
facilitated model out of equilibrium, \emph{Markov Process.\ Related Fields} \textbf{19} (2013), no.~3, 383--406.
\arxiv{1205.4584}
\bibitem{BDT19} O.~Blondel, A.~Deshayes and C.~Toninelli, Front evolution of the Fredrickson--Andersen one spin
facilitated model, \emph{Electron.\ J.\ Probab.} \textbf{24} (2019), paper no.~1, 32 pp. \doi{10.1214/18-EJP246}
\bibitem{CMRT08} N.~Cancrini, F.~Martinelli, C.~Roberto and C.~Toninelli, Kinetically constrained spin models,
\emph{Probab.\ Theory Related Fields} \textbf{140} (2008), no.~3--4, 459--504. \doi{10.1007/s00440-007-0072-3}
\bibitem{CMST10} N.~Cancrini, F.~Martinelli, R.~Schonmann and C.~Toninelli, Facilitated oriented spin models: some
non equilibrium results, \emph{J.\ Stat.\ Phys.} \textbf{138} (2010), no.~6, 1109--1123.
\doi{10.1007/s10955-010-9923-x}
\bibitem{FA84} G.~H.~Fredrickson and H.~C.~Andersen, Kinetic Ising model of the glass transition, \emph{Phys.\ Rev.\
Lett.} \textbf{53} (1984), no.~13, 1244--1247. \doi{10.1103/PhysRevLett.53.1244}
\bibitem{Mar26} F.~Martinelli, Kinetically constrained spin models: a short survey, in \emph{Proceedings of the
International Congress of Mathematicians 2026}, vol.~5, 497--515. \doi{10.1137/25M180425X}
\bibitem{MST25} F.~Martinelli, A.~Shapira and C.~Toninelli, Long time behaviour of one facilitated kinetically
constrained models: results and open problems, preprint, 2025. \arxiv{2510.20461}
\bibitem{MV19} T.~Mountford and G.~Valle, Exponential convergence for the Fredrickson--Andersen one-spin facilitated
model, \emph{J.\ Theoret.\ Probab.} \textbf{32} (2019), no.~1, 282--302. \doi{10.1007/s10959-018-0812-6}
\bibitem{MathBrainRepo} Y.~Shmalo, \emph{Math Brain: public snapshot, code and tests}, 2026,
\url{https://github.com/yspennstate/math-brain-public}.
\bibitem{MathBrainPaper} Y.~Shmalo, \emph{Math Brain: shared, auditable memory for multi-agent mathematical research},
October 2026, \url{https://github.com/yspennstate/math-brain-public/tree/main/papers/math-brain-2026-10-24}.
\bibitem{BrainFramework} Y.~Shmalo, \emph{Brain Framework}, public software repository, 2026,
\url{https://github.com/yspennstate/brain-framework-public}.
\end{thebibliography}
\end{document}
